\documentclass{beamer}
\usepackage{tikz,amsmath,hyperref,graphicx,stackrel,animate,amssymb}
\usetikzlibrary{positioning,shadows,arrows,shapes,calc,dsp}
\newcommand{\argmax}{\operatornamewithlimits{argmax}}
\newcommand{\argmin}{\operatornamewithlimits{argmin}}
\mode<presentation>{\usetheme{Frankfurt}}
\AtBeginSection[]
{
  \begin{frame}<beamer>
    \frametitle{Outline}
    \tableofcontents[currentsection,currentsubsection]
  \end{frame}
}
\title{Lecture 25: Exam 2 Review}
\author{Mark Hasegawa-Johnson}
\date{ECE 401: Signal and Image Analysis}
\begin{document}

% Title
\begin{frame}
  \maketitle
\end{frame}

% Title
\begin{frame}
  \tableofcontents
\end{frame}

%%%%%%%%%%%%%%%%%%%%%%%%%%%%%%%%%%%%%%%%%%%%
\section[Administration]{Exam administration}
\setcounter{subsection}{1}

\begin{frame}
  \frametitle{When, Where}

  \begin{itemize}
  \item When: Monday, November 2, 1:00-1:50pm.
  \item Where: ECEB 3015.
  \item Online section: Be signed on to zoom 12:55pm, with your camera
    on.  Have Gradescope open on your phone or device, so you can
    upload your answers promptly at 1:50pm when the exam is completed.
  \end{itemize}
\end{frame}

\begin{frame}
  \frametitle{What to bring}

  \begin{itemize}
  \item Bring: Pencil or pen.  Erasers.
  \item Notes sheet: You'll be given one formula sheet.  You are also permitted
    to bring your own formula sheet.  Yours should be 8.5x11, handwritten, both sides.
  \item Don't bring: Anything electronic.  Any of your own paper (the
    last page of the exam paper will be blank, so you can use it as
    scratch paper).
  \end{itemize}
\end{frame}

\begin{frame}
  \frametitle{Grading}

  \begin{itemize}
  \item Exam will be graded on Gradescope.
  \item There will probably be four problems, each worth about 25 points.
  \item Grades should be available one week after the exam.
  \end{itemize}
\end{frame}

%%%%%%%%%%%%%%%%%%%%%%%%%%%%%%%%%%%%%%%%%%%%
\section[Topics]{Topics included in the exam}
\setcounter{subsection}{1}

\begin{frame}
  \frametitle{Topics included in the exam: Lectures 15-22}

  \begin{enumerate}
  \item DTFT, Windowing, and Windows
  \item Filter design: Ideal filters, Windowing
  \item DTFT of tones: Delta function, Fourier series, Spectral estimation
  \item Overlap-Add
  \end{enumerate}
\end{frame}

\begin{frame}
  \frametitle{DTFT}

  The DTFT (discrete time Fourier transform) of any signal is
  $X(\omega)$, given by
  \begin{align*}
    X(\omega) &= \sum_{n=-\infty}^\infty x[n]e^{-j\omega n}\\
    x[n] &= \frac{1}{2\pi}\int_{-\pi}^\pi X(\omega)e^{j\omega n}d\omega
  \end{align*}
  Particular useful examples include:
  \begin{align*}
    f[n]=\delta[n] &\leftrightarrow F(\omega)=1\\
    g[n]=\delta[n-n_0] &\leftrightarrow G(\omega)=e^{-j\omega n_0}
  \end{align*}
\end{frame}

\begin{frame}
  \frametitle{Properties of the DTFT}

  Properties worth knowing  include:
  \begin{enumerate}
    \setcounter{enumi}{-1}
  \item Periodicity: $X(\omega+2\pi)=X(\omega)$
  \item Linearity:
    \[z[n]=ax[n]+by[n]\leftrightarrow Z(\omega)=aX(\omega)+bY(\omega)
    \]
  \item Time Shift: $x[n-n_0]\leftrightarrow e^{-j\omega n_0}X(\omega)$
  \item Frequency Shift: $e^{j\omega_0 n}x[n]\leftrightarrow X(\omega-\omega_0)$
  \item Filtering and Windowing:
    \begin{align*}
    y[n]=h[n]\ast x[n]&\leftrightarrow Y(\omega)=H(\omega)X(\omega)\\
    y[n]=w[n]x[w] &\leftrightarrow Y(\omega)=\frac{1}{2\pi}W(\omega)\ast X(\omega)
    \end{align*}
  \end{enumerate}
\end{frame}

\begin{frame}
  \frametitle{Ideal Filters}
  \begin{itemize}
  \item Ideal Lowpass Filter:
    \[
    H_{LP}(\omega)
    = \begin{cases} 1& |\omega|\le\omega_c,\\
      0 & \omega_c<|\omega|\le\pi.
    \end{cases}~~~\leftrightarrow~~~
    h_{LP}[m]=\frac{\omega_c}{\pi}\mbox{sinc}(\omega_c n)
    \]
  \item Ideal Highpass Filter:
    \[
    H_{HP}(\omega)=1-H_{LP}(\omega)~~~\leftrightarrow~~~
    h_{HP}[n]=\delta[n]-\frac{\omega_c}{\pi}\mbox{sinc}(\omega_c n)
    \]
  \item Ideal Bandpass Filter:
    \begin{align*}
      H_{BP}(\omega)&=H_{LP,\omega_2}(\omega)-H_{LP,\omega_1}(\omega)\\
      \leftrightarrow
      &h_{BP}[n]=\frac{\omega_2}{\pi}\mbox{sinc}(\omega_2 n)-\frac{\omega_1}{\pi}\mbox{sinc}(\omega_1 n)
    \end{align*}
  \end{itemize}
\end{frame}

\begin{frame}
  \frametitle{Windows}
  \begin{itemize}
  \item Rectangular Window:
    \[
    w_R[n] = \begin{cases}
      1 & -\left(\frac{L-1}{2}\right)\le n\le \left(\frac{L-1}{2}\right)\\
      0 & \mbox{otherwise}
    \end{cases}~~~\leftrightarrow~~~
    W_R(\omega)=\frac{\sin(\omega L/2)}{\sin(\omega/2)}
    \]
  \item Hamming and Hann Windows:
    \[
    w_H[n] = w_R[n]\left(A+(1-A)\cos\left(\frac{2\pi n}{L-1}\right)\right)
    \]
    \begin{align*}
    \leftrightarrow~~~
    W_H(\omega)&=AW_R(\omega)+\frac{1-A}{2}W_R\left(\omega-\frac{2\pi}{L-1}\right)\\
    &+\frac{1-A}{2}W_R\left(\omega+\frac{2\pi}{L-1}\right)
    \end{align*}
  \end{itemize}
\end{frame}

\begin{frame}
  \frametitle{Pure Tones}

  \begin{itemize}
  \item DTFT of a complex exponential is a delta function:
    \begin{displaymath}
      e^{j\omega_0 n}~~\leftrightarrow~~~2\pi\delta(\omega-\omega_0)
    \end{displaymath}
  \item DTFT of a cosine is two delta functions:
    \begin{displaymath}
      \cos(\omega_0 n)~~\leftrightarrow~~~\pi\delta(\omega-\omega_0)+\pi\delta(\omega+\omega_0)
    \end{displaymath}
  \item DTFT of a windowed cosine is frequency-shifted window functions:
    \begin{displaymath}
      \cos(\omega_0 n)w[n]~~\leftrightarrow~~~\frac{1}{2}W(\omega-\omega_0)+\frac{1}{2}W(\omega+\omega_0)
    \end{displaymath}
  \end{itemize}
\end{frame}

\begin{frame}
  \frametitle{Complex Tones}
  If $x[n]$ is periodic with Fourier series coefficients, then its DTFT and windowed DTFT are
  \begin{align*}
    \mathcal{F}\left\{x[n]\right\}
    &=2\pi\sum_{k=-N_0/2}^{(N_0-1)/2} X_k \delta\left(\omega-\frac{2k\pi}{N_0}\right)\\
    \mathcal{F}\left\{w[n]x[n]\right\}
    &=\sum_{k=-N_0/2}^{(N_0-1)/2} X_k W\left(\omega-\frac{2k\pi}{N_0}\right)
  \end{align*}
\end{frame}

\begin{frame}
  \frametitle{The Overlap-Add Algorithm}

  \begin{enumerate}
  \item Divide $x[n]$ into frames of length $M$, zero-pad to $N\ge L+M-1$, then FFT:
    \begin{align*}
      x_t[n] &= x[n+tM]w[n]\\
      X_t[k] &= \text{FFT}\{x_t[n]\}
    \end{align*}
  \item Generate the output from each frame
    \begin{align*}
      Y_t[k] &= X_t[k]H[k]\\
      y_t[n] &= \text{FFT}^{-1}\{y_t[n]\}
    \end{align*}
  \item  Overlap the outputs, and add them together
    \begin{align*}
      y[n] &= \sum_t y_t[n-tM]
    \end{align*}
  \end{enumerate}
\end{frame}

%%%%%%%%%%%%%%%%%%%%%%%%%%%%%%%%%%%%%%%%%%%%
\section[Sample problems]{Sample problems}
\setcounter{subsection}{1}

\begin{frame}
  \frametitle{Sample problems}

  If there's time, at this point in the lecture, I will do some sample
  problems from the practice exam.  If there are problems whose
  solution you'd like to see, please ask.
\end{frame}

\end{document}
